% HiMCM 2020 Problem A --- two-page student guide (Requirement 5)
% Non-technical. True ranks only. ~80% ASD-STE100.
\documentclass[11pt]{article}
\usepackage[margin=0.85in]{geometry}
\usepackage{booktabs}
\usepackage{enumitem}
\usepackage[hidelinks]{hyperref}
\pagestyle{myheadings}
\markright{Team XXXX \hfill A Data-Driven Summer Job Guide}

\begin{document}
\begin{center}
{\Large\bfseries Navigating Your First Paycheck:}\\[0.2em]
{\large A Data-Driven Guide to Finding the Ideal Summer Job}\\[0.4em]
{\small HiMCM 2020 Problem A \quad Team XXXX}
\end{center}

\medskip
\noindent
This note is for students, not for a statistics class.
We ranked eight common summer job candidates with a small data pipeline.
The pipeline is imperfect.
Read the three traps first.
Then use the web tool in three steps.

\section*{What the data actually said}

We used 50 student rows and eight job ads (lifeguard, camp counselor, tutor, fast-food crew, retail clerk, caddy, pet sitter, swim coach).
Pay snippets come from those ads.
Example: a municipal lifeguard ad lists about \(\$15.50\) to \(\$18.00\) per hour.
A tutor ad lists about \(\$22\) to \(\$35\) per hour.

A simple score that rewards skill/network, rewards schedule control, and penalizes physical load put \textbf{tutor} first and \textbf{lifeguard} last.
That score is not ``who you should be.''
It is one ranking under one formula.

\section*{Three traps}

\textbf{Trap 1 --- Pretty names can be wrong.}
A slide deck may call three scores ``money,'' ``resume,'' and ``pain.''
Our rotation did not recover that split.
Safety, friends, and skills sat on the same score.
Hourly wage sat with commute and control, not with tips alone.
If a poster says ``Factor 1 \(=\) cash,'' ask to see the loading table.

\textbf{Trap 2 --- A high hourly rate is not the whole job.}
Caddie loops can pay a lot on a busy Saturday.
The same loop is sun, hills, and a heavy bag.
Fast food is easier to get and harder on your feet.
Tutor pay is high and the work is mostly indoor and seated.
Write down cash, hours, heat, and what you can put on an application.

\textbf{Trap 3 --- A computer rank is not a hiring officer.}
Our classifier saw 50 labeled students.
Its best-guess hit rate for the single true job was \(20\%\).
Guessing the most common job (retail) hit \(22\%\).
So the model is weaker than a simple ``most people picked retail'' rule.
Use it to compare options.
Do not treat a \(40\%\) match as a promise.

\section*{Three-step web tool}

The tool is \texttt{SummerJobMatch}.
On a contest laptop run:
\texttt{streamlit run src/web\_app/app.py}.
Then:

\begin{enumerate}[leftmargin=1.4em,itemsep=0.35em]
\item \textbf{Pick a preset or move three sliders (1 to 10).}
``I need money'' sets slider 1 to 9 and slider 3 to 8.
``I need application stories'' sets slider 2 to 10.
``I want an easy summer'' sets slider 3 to 1.
\item \textbf{Read three cards.}
Gold \(=\) best match under the model.
Silver \(=\) runner-up.
Red \(=\) the strongest mismatch.
Each card shows a pay span copied from the job ad.
\item \textbf{Ignore the nickname if the card fights your body.}
Example: the ``easy summer'' preset still ranked \textbf{lifeguard} first in our run (\(40.9\%\)).
That is a model glitch, not advice to sit in the sun all day.
If you hate heat, throw that rank out.
\end{enumerate}

\section*{A short checklist before you apply}

\begin{itemize}[leftmargin=1.2em,itemsep=0.25em]
\item Can you name the cash range from a real posting, not a rumor?
\item Can you work the posted hours, including weekends?
\item What skill or story will you write on a college form in September?
\item What is the worst physical day in that job, and have you tried a sample of it?
\end{itemize}

\noindent
If those four answers are clear, you do not need a neural net.
The tool is a map.
You still choose the road.

\end{document}
